% =====================================================================
%  RESUMEN Y ABSTRACT
%  La plantilla fija 150-250 palabras y pide objetivo, método, resultados y
%  discusión. Las palabras clave son obligatorias en los dos idiomas.
% =====================================================================
\ustaseccion{Resumen}

Orientarse en un edificio de varios pisos es difícil para quien no lo conoce, y
el cambio de nivel corta la continuidad del recorrido. Las soluciones robóticas
existentes resuelven ese salto dotando al robot de mecanismos para subir
escaleras o conectándolo al ascensor, dos vías caras que obligan a intervenir la
infraestructura. Este trabajo propone y evalúa una alternativa: un sistema
colaborativo de dos robots móviles terrestres, cada uno dedicado a una planta,
que se reparten el guiado de una persona mediante un protocolo de relevo en la
zona de escaleras. Ningún vehículo cruza de nivel. La arquitectura es
centralizada: un coordinador asigna el agente según el nivel del origen y del
destino, y el usuario solicita el servicio desde una interfaz web que funciona en
cualquier teléfono sin instalación ni acceso a internet. La evaluación siguió un
protocolo escrito antes de instrumentar el sistema, con cuatro métricas de
definición operativa y treinta misiones sorteadas en simulación. La campaña
resultó válida sin descartes y permite afirmar, con 95~\% de confianza, que la
tasa de éxito supera el 70~\%. Sobre los vehículos físicos se verificó la cadena
completa de navegación y la odometría cumplió su criterio, con error entre
$-4{,}6$~\% y $+3{,}3$~\% sobre recorridos de 6,7 y 14,6~m. La precisión de
llegada quedó fuera de la tolerancia por una limitación medida de la plataforma:
el vehículo no admite velocidades por debajo de 0,40~m/s y carece de régimen de
aproximación fina.

\vspace{1em}
\textbf{Palabras clave:} robótica móvil colaborativa; sistemas multi-robot;
navegación en interiores; protocolo de relevo; interacción humano-robot; ROS~2.

\ustaseccion{Abstract}

Finding one's way inside a multi-storey building is hard for first-time
visitors, and changing floors breaks the continuity of the route. Existing
robotic solutions address that gap either by equipping the robot to climb stairs
or by interfacing it with the building lift, two expensive approaches that
require modifying the infrastructure. This work proposes and evaluates an
alternative: a collaborative system of two ground mobile robots, each assigned to
one floor, that share the task of guiding a person through a handover protocol at
the stairwell. Neither vehicle crosses between levels. The architecture is
centralised: a coordinator assigns the agent according to the floor of the origin
and the destination, and the user requests the service from a web interface that
runs on any phone with no installation and no internet access. The evaluation
followed a protocol written before instrumenting the system, with four
operationally defined metrics and thirty randomised missions in simulation. The
campaign was valid with no discards and supports, with 95~\% confidence, the
claim that the success rate exceeds 70~\%. On the physical vehicles the full
navigation chain was verified and the odometry met its criterion, with errors
between $-4{,}6$~\% and $+3{,}3$~\% over runs of 6.7 and 14.6~m. Arrival accuracy
fell outside tolerance because of a measured platform limitation: the vehicle
does not accept commanded speeds below 0.40~m/s and therefore has no fine
approach regime.

\vspace{1em}
\textbf{Keywords:} collaborative mobile robotics; multi-robot systems; indoor
navigation; handover protocol; human-robot interaction; ROS~2.
